% TOP_PROBLEM: {"id": 30007141, "title": "Unique multiplicity-one tangents at first surface-flow singularities", "category_id": 6, "category": {"id": 6, "name": "geometry", "display_name": "Geometry", "description": "Euclidean and non-Euclidean geometry, geometric structures.", "slug": "geometry", "order_index": 6, "created_at": "2026-07-31T15:26:25.671Z"}, "status": "open"}
% FIRST_POSTED: 2026-10-06
% ATTEMPT_STATUS: solved
% OPENAI_CLAIM: {"schema_version":1,"source_repo":"https://github.com/openai/math","source_commit":"adc7f1241b42e322a6451854ab7e4b4c146bf78a","release_date":"2026-10-06","match":"full","resolution":"proved","independently_reviewed":false,"second_pass":"agreed","families":[{"family":"355","title":"Unique tangent flows at the first surface singularity","subject":"Differential geometry","manuscripts":[{"title":"Uniqueness of tangent flows at the first singular time of embedded surface mean-curvature flow","pdf_path":"preprints/Tangent-flow-uniqueness-2026-09-24/paper.pdf","readme_path":"preprints/Tangent-flow-uniqueness-2026-09-24/README.md","date":"2026-09-24","theorem_ref":"thm:main (sections/introduction.tex:46)","statement_sources":[{"path":"preprints/Tangent-flow-uniqueness-2026-09-24/build/sections/introduction.tex","line":46,"label":"thm:main"}]}],"lean":null,"reasoning_trace":null}]}
\documentclass[11pt]{article}
\usepackage[margin=25mm]{geometry}
\usepackage{fontspec}
\setmainfont{Latin Modern Roman}
\usepackage{amsmath,amssymb,mathtools}
\usepackage{unicode-math}
\setmathfont{Latin Modern Math}
\usepackage{enumitem}
\usepackage{xurl,hyperref,bookmark}
\hypersetup{colorlinks=true,urlcolor=blue}
\setcounter{secnumdepth}{0}
\setlength{\parindent}{0pt}
\setlength{\parskip}{5pt}
\setlength{\emergencystretch}{3em}
\begin{document}
\section*{Unique multiplicity-one tangents at first surface-flow singularities}

\subsection{OpenAI's claim}
\textit{Uniqueness of surface tangent flow}: At the first singular time of a smooth closed embedded surface flow in $\mathbb R^3$, the paper claims a unique multiplicity-one self-shrinking tangent flow at each fixed singular point. Its position and axes are unique in the original coordinates.

\subsection{Mathematical statement of the claimed result}
\paragraph{Mathematical claim 1.} Let a smooth compact connected embedded boundaryless surface $M_0\subset\mathbb R^3$ evolve by its maximal smooth mean-curvature flow to first singular time $T$. At every singular $(x_0,T)$, the claim provides one smooth properly embedded self-shrinker $S$ such that \[\lambda(M_{T+s/\lambda^2}-x_0)\longrightarrow\sqrt{-s}\,S\quad(\lambda\to\infty,\ s<0).\] Convergence holds for the full dilation family, locally smoothly in space with multiplicity one and uniformly on compact negative-time intervals, in the original ambient coordinates. The backward Brakke tangents quantified here are weakly measurable Radon families $(\mu_s)_{s<0}$ obtained as spacetime weight limits of dilation subsequences, with integral two-varifold slices for almost every $s$, generalized mean curvature $\mathbf H_s$, and $\delta V_s=-\mathbf H_s\mu_s$. For each compact spatial $K$ and compact time interval $J\subset(-\infty,0)$ they obey $\int_J\int_K(1+|\mathbf H_s|^2)\,d\mu_s\,ds<\infty$. Every choice of $a<b<0$ and nonnegative $\varphi\in C_c^1(\mathbb R^3\times[a,b])$ obeys \[\mu_b(\varphi_b)-\mu_a(\varphi_a)\le\int_a^b\int(\partial_s\varphi+\nabla\varphi\cdot\mathbf H_s-\varphi|\mathbf H_s|^2)\,d\mu_s\,ds.\] Tests may be nonzero at $a,b$, and these endpoints need not be typical slice times. All such tangents satisfy $\mu_s(A)=\mathcal H^2(A\cap\sqrt{-s}\,S)$ for every Borel $A$ and every $s<0$. Later singularities, higher dimensions and arbitrary weak initial flows are not included.

\subsection{Outline of the argument}
\textit{Uniqueness of surface tangent flow}: The stationary analysis handles conical and cylindrical shrinker ends together, including quadratic opening modes for which ordinary analytic graph arguments do not suffice. Finite Taylor jets and exponentially small end terms produce a localized gradient inequality. Increasing graph radii are chosen from preceding energy drops so that cutoff errors are summable; finite rescaled length gives convergence with its ambient position and axes fixed.

\subsection{Scope relative to this problem}
\textit{Uniqueness of surface tangent flow}: At the first singular time of a smooth closed embedded surface flow in $\mathbb R^3$, the paper claims a unique multiplicity-one self-shrinking tangent flow at each fixed singular point. Its position and axes are unique in the original coordinates. The parameters, regularity assumptions and normalization in this question are retained in applying that assertion.

With the hypotheses and conventions above, the claimed result addresses the whole mathematical assertion on this page.

\subsection{Formal verification}
\textit{Uniqueness of surface tangent flow}: The pinned catalogue lists no Lean formalization for this result.

This site has not run the Lean development. These links describe OpenAI’s selected formal statements and should be read with their stated scope.

\subsection{Status caveat}
OpenAI reports that an internal model produced the manuscripts. Their arguments have not been independently verified as proofs by this site. Any portion outside the documented formal scope may contain mistakes, and the catalogue label reports the announced mathematical claim.

\subsection{Sources}
\begin{itemize}
\item OpenAI, \emph{Uniqueness of tangent flows at the first singular time of embedded surface mean-curvature flow}, 2026-09-24; citation key \texttt{OAI:Tangent-flow-uniqueness-2026-09-24}. \href{https://github.com/openai/math/blob/adc7f1241b42e322a6451854ab7e4b4c146bf78a/preprints/Tangent-flow-uniqueness-2026-09-24/paper.pdf}{PDF (GitHub)}; \href{https://github.com/openai/math/raw/adc7f1241b42e322a6451854ab7e4b4c146bf78a/preprints/Tangent-flow-uniqueness-2026-09-24/paper.pdf}{PDF (direct)}; \href{https://github.com/openai/math/blob/adc7f1241b42e322a6451854ab7e4b4c146bf78a/preprints/Tangent-flow-uniqueness-2026-09-24/README.md}{Manuscript page and citation}; \href{https://github.com/openai/math/blob/main/preprints/Tangent-flow-uniqueness-2026-09-24/paper.pdf}{Latest version}; Statement source: \href{https://github.com/openai/math/blob/adc7f1241b42e322a6451854ab7e4b4c146bf78a/preprints/Tangent-flow-uniqueness-2026-09-24/build/sections/introduction.tex\#L46}{thm:main}.
\item \href{https://github.com/openai/math/blob/adc7f1241b42e322a6451854ab7e4b4c146bf78a/CONTENTS.md}{OpenAI catalogue entry}.
\end{itemize}
\end{document}
